% GENERATED FILE. Edit canonical lessons, then run npm run generate:books.
% canonical-source-hash: fnv1a64:2bd71dacbad3e98e
% canonical-lessons: MR-C68-shubh-sakal, MR-C68-shubh-dupar, MR-C68-shubh-sandhyakal, MR-R68-greetings-recall

\chapter{Three Greetings, One Word}
\label{ch:mr-shubh}
\hlchaptermodality{68}

\begin{chapteropening}
\textbf{By the end of this chapter:} \emph{I can read the time-of-day greetings, and I know which greeting Marathi speakers actually use at any hour.}

From cold: the three shubh greetings in the order the day runs, the one new word they are all built from, and the greeting that replaces all of them in speech.
\end{chapteropening}

\section[śubh sakāḷ]{\mr{शुभ} \mr{सकाळ} (śubh sakāḷ) — \textquotedblleft{}good morning\textquotedblright{}}

\label{lesson:MR-C68-shubh-sakal}



\begin{quote}
Say the word for morning. Then say the greeting you have used since your very first Marathi word.
\end{quote}

\subsection*{You'll want to know: \mr{शुभ}}
\begin{quote}
\textbf{\mr{शुभ}} — \emph{śubh} — \textbf{auspicious, good}
\end{quote}

One word. It goes in front of a part of the day, and that is the whole construction:

\begin{quote}
\textbf{\mr{शुभ} \mr{सकाळ}.} — \emph{śubh sakāḷ.} — \textbf{good morning.}
\end{quote}

The second word is not new. \textbf{\mr{सकाळ}} is yours already, and so are the other parts of the day — which means \textbf{one new word buys you a greeting for each of them.} That is the cheapest chapter you have had in a while, and it is cheap because of work you already did.

\begin{culture}
Here is the part worth having, and it is about \textbf{when you would actually say this.}

\textbf{\mr{नमस्कार} covers every hour of the day.} Your first Marathi word was the all-purpose greeting, and it has not stopped being that. Marathi does not need a different greeting for the morning the way English does.

\textbf{\mr{शुभ} \mr{सकाळ}} is real, and it skews \textbf{formal and written} — a message, a notice, a broadcast, the top of something addressed to a group. Spoken between two people who know each other, \textbf{\mr{नमस्कार}} is the ordinary choice at any hour.

So learn this to \textbf{read} it and to use it where it belongs, and do not replace the greeting you already have.
\end{culture}

\subsection*{Your turn}
Say these aloud:

\begin{itemize}
\raggedright
  \item \emph{śubh} — good, auspicious
  \item \emph{śubh sakāḷ}
  \item which of the two words in it you already knew
  \item what you would say instead, out loud, to a person
\end{itemize}

\begin{tcolorbox}[breakable,colback=teal!4,colframe=teal!35!black,title={Before you move on}]
What does \textbf{\mr{शुभ}} mean on its own? (\textbf{Auspicious, good}.) How many new words did \emph{good morning} cost you? (\textbf{One}.) And which greeting still covers every hour? (\emph{\textbf{\mr{नमस्कार}}}.)
\end{tcolorbox}
\section[śubh dupār]{\mr{शुभ} \mr{दुपार} (śubh dupār) — \textquotedblleft{}good afternoon\textquotedblright{}}

\label{lesson:MR-C68-shubh-dupar}



\begin{quote}
Say \emph{good morning}. Now say the part of the day that comes after the morning, and put the same word in front of it before you read on.
\end{quote}

\subsection*{You'll want to know: \mr{शुभ} \mr{दुपार}}
\begin{quote}
\textbf{\mr{शुभ} \mr{दुपार}.} — \emph{śubh dupār.} — \textbf{good afternoon.}
\end{quote}

If you built that yourself in the warm-up, the lesson has already done its job. \textbf{\mr{शुभ}} does not change for the word behind it, and the word behind it is one you have had since the parts of the day.

\noindent
\begin{tabularx}{\linewidth}{@{}>{\raggedright\arraybackslash}X>{\raggedright\arraybackslash}X@{}}
\toprule
\textbf{} & \textbf{} \\
\midrule
\textbf{\mr{शुभ} \mr{सकाळ}} & good morning \\
\textbf{\mr{शुभ} \mr{दुपार}} & good afternoon \\
\bottomrule
\end{tabularx}

\begin{grammarlens}[title={Grammar Lens: a pattern you can finish alone}]
Two rows are enough to see the shape: \textbf{\mr{शुभ}} plus a part of the day.

\textbf{You can already say the evening greeting}, and the next lesson will only confirm what you worked out rather than teach it to you. That is what a system is worth — the words stop arriving one at a time.

This one is also the \textbf{least used} of the set in speech, for the same reason the last lesson gave: there is a greeting that covers the whole day and Marathi speakers reach for it.
\end{grammarlens}

\subsection*{Your turn}
Say these aloud:

\begin{itemize}
\raggedright
  \item \emph{śubh dupār}
  \item both greetings you have, in the order the day runs
  \item the third one, before anybody teaches it to you
\end{itemize}

\begin{tcolorbox}[breakable,colback=teal!4,colframe=teal!35!black,title={Before you move on}]
How do you say \emph{good afternoon}? (\emph{\textbf{\mr{शुभ} \mr{दुपार}}}.) Did \textbf{\mr{शुभ}} change? (\textbf{No} — it never does.) Which part of the day can you still put behind it? (\textbf{The evening one}.)
\end{tcolorbox}
\section[śubh sandhyākāḷ]{\mr{शुभ} \mr{संध्याकाळ} (śubh sandhyākāḷ) — \textquotedblleft{}good evening\textquotedblright{}}

\label{lesson:MR-C68-shubh-sandhyakal}



\begin{quote}
Say the two greetings you have. Then say the one you worked out last lesson without being taught it.
\end{quote}

\subsection*{You'll want to know: \mr{शुभ} \mr{संध्याकाळ}}
\begin{quote}
\textbf{\mr{शुभ} \mr{संध्याकाळ}.}
\end{quote}

\begin{quote}
— \emph{śubh sandhyākāḷ.} — \textbf{good evening.}
\end{quote}

You built this before you read it. \textbf{\mr{शुभ}} in front, the part of the day behind, nothing changed in either.

\noindent
\begin{tabularx}{\linewidth}{@{}>{\raggedright\arraybackslash}X>{\raggedright\arraybackslash}X@{}}
\toprule
\textbf{} & \textbf{} \\
\midrule
\textbf{\mr{शुभ} \mr{सकाळ}} & good morning \\
\textbf{\mr{शुभ} \mr{दुपार}} & good afternoon \\
\textbf{\mr{शुभ} \mr{संध्याकाळ}} & good evening \\
\bottomrule
\end{tabularx}

\begin{culture}
Three greetings, and the honest summary of all three is the same.

They belong to \textbf{writing and to formal speech} — a notice, a message, the opening of something addressed to a room. Between two people, at any hour of the day, the greeting is \textbf{\mr{नमस्कार}}, and it was the first word this book gave you.

That is not a reason to skip them. \textbf{You will read them}, and a reader who has never seen \textbf{\mr{शुभ}} will stall on a word that costs nothing to know. It is a reason not to swap out a working greeting for one that will sound like a broadcast.
\end{culture}

\begin{grammarlens}[title={Grammar Lens: where the pattern stops}]
You have four parts of the day and three greetings. \textbf{The night is the one that does not follow.}

Its greeting is built on \textbf{\mr{रात्री}} — the at-that-time form from the parts of the day — rather than on the bare \textbf{\mr{रात्र}} these three use. This book is not teaching that phrase here, and it is naming the exception so that you do not build it yourself and get it wrong.

\textbf{A pattern worth trusting is one you also know the edge of.}
\end{grammarlens}

\subsection*{Your turn}
Say these aloud:

\begin{itemize}
\raggedright
  \item \emph{śubh sandhyākāḷ}
  \item all three, in the order the day runs
  \item which part of the day has no greeting in this chapter
  \item what you would actually say to a person, at any hour
\end{itemize}

\begin{tcolorbox}[breakable,colback=teal!4,colframe=teal!35!black,title={Before you move on}]
Say \emph{good evening}. (\emph{\textbf{\mr{शुभ} \mr{संध्याकाळ}}}.) Which register do these three belong to? (\textbf{Written and formal}.) And which part of the day was left out, because it does not follow the pattern? (\textbf{The night}.)
\end{tcolorbox}
\section[śubh ābhivādan]{(three greetings, one word) — from cold}

\label{lesson:MR-R68-greetings-recall}



\begin{quote}
Close the lessons before this one. Say the one word this chapter taught, and then say what it means on its own.
\end{quote}

\begin{grammarlens}[title={Grammar Lens: one word, three greetings}]
\noindent
\begin{tabularx}{\linewidth}{@{}>{\raggedright\arraybackslash}X>{\raggedright\arraybackslash}X@{}}
\toprule
\textbf{} & \textbf{} \\
\midrule
\textbf{\mr{शुभ} \mr{सकाळ}} & good morning \\
\textbf{\mr{शुभ} \mr{दुपार}} & good afternoon \\
\textbf{\mr{शुभ} \mr{संध्याकाळ}} & good evening \\
\bottomrule
\end{tabularx}

\textbf{This chapter taught one word.} Everything on the right of \textbf{\mr{शुभ}} you already had, which is why three greetings cost what one usually does.

That is the second time the parts of the day have paid for something: they gave you the at-that-time forms first, and now they give you these.
\end{grammarlens}

\begin{grammarlens}[title={What you've built}]
Three greetings \textbf{to read}, and one greeting \textbf{to say}.

\textbf{\mr{नमस्कार}} was the first Marathi word in this book and it still covers every hour of the day. The three above belong to notices, messages and formal openings. Knowing both, and knowing which is which, is the whole of what this chapter is for.

The night has no greeting here on purpose: its form is built on \textbf{\mr{रात्री}} rather than the bare part of the day, so it does not follow these three.
\end{grammarlens}

\subsection*{Your turn}
Say these aloud:

\begin{itemize}
\raggedright
  \item the three greetings, in the order the day runs
  \item the one word that is new, and what it means alone
  \item which greeting you would use out loud, and at what hour
  \item which part of the day is missing from the three, and why
\end{itemize}

\begin{tcolorbox}[breakable,colback=teal!4,colframe=teal!35!black,title={Before you move on}]
How many new words did this chapter teach? (\textbf{One} — \textbf{\mr{शुभ}}.) Say \emph{good evening}. (\emph{\textbf{\mr{शुभ} \mr{संध्याकाळ}}}.) And which greeting would you actually speak, at any hour of the day? (\emph{\textbf{\mr{नमस्कार}}}.)
\end{tcolorbox}
